% Nouveau EGA V (cote 25): archive cote 25, site page 182 (PDF page 183).
% English translation of the checked French annotated master.
% AI-translated by GPT-6 Astra (OpenAI, ChatGPT Pro mode).
\scanpage{25/182}
\mspage{\hand{9}}

\unsure{accepted}\textsubscript{\unsure{,}} terminology---I do not have Bourbaki at hand to check this
point---) should be called the \uline{determinant} of the quadratic form, and not
the discriminant. I should like to know your opinion on this question; if
you agree, \unsure{w}\unsure{e} would take this opportunity to correct the
accepted terminology on this point, since it is misleading (because until
the last few days I had unknowingly confused discriminant and
determinant).

\uline{Remarks} 1\var{\over{\unsure{5}}{6}}{6}{writing over}{ega25_25_182-v0001}.12. When one lets the absolute group scheme
$Gl(n)$ \var{\deleted{\var{\xout{over Z}}{}{typewriter strike-out}{ega25_25_182-v0003}}}{}{deletion}{ega25_25_182-v0002} $= \underline{Aut}(\underline{Z}^n )$ act on $\underline{Q}(\underline{Z}^n)$ , the stabilizer of the standard quadratic
form is denoted by $\underline{O}(n)$, and is called the ``\uline{absolute orthogonal}
\uline{group}''.( More generally, for every quadratic form $Q$ \hand{$\in$}
$\var{\over{T}{\unsure{\Gamma}}}{\unsure{\Gamma}}{writing over}{ega25_25_182-v0004}$ $\underline{Sym}^2(\underline{E})$ , one introduces the closed group subprescheme of
$Gl(\underline{E})$ \var{\xout{\unsure{which}}}{}{typewriter strike-out}{ega25_25_182-v0005} stabilizing $Q$ , denoted by $\underline{O}(Q)$ and called the \uline{orthogonal group scheme}
\uline{relative to the form Q}; when $Q$ is isomorphic to the inverse
image of the standard form \var{\inserted[typed]{in n variables}}{in n variables}{insertion}{ega25_25_182-v0006}\hand{,} $\underline{O}(Q)$ is isomorphic to $O(n)_S$ ). This being so,
it is immediate that two points \var{\inserted[typed]{$\unsure{Q|u},\unsure{Q|u}'$}}{$\unsure{Q|u},\unsure{Q|u}'$}{insertion}{ega25_25_182-v0007} of the first member of (**), with values in
an $S'$, have the same image in the second member iff \var{\xout{they are congru\unsure{ent} \unsure{the}}}{}{typewriter strike-out}{ega25_25_182-v0008}
\var{\xout{second}}{}{typewriter strike-out}{ega25_25_182-v0009} we have $u' = uv$ , with $v \hand{\in} \underline{O}(n)(S')$; in view of 16.11. (and
with the terminology that will be developed in greater detail in
EGA V or VI\var{\inserted{\unsure{(}see in the meantime SGAD \uline{IV}\unsure{)}}}{\unsure{(}see in the meantime SGAD \uline{IV}\unsure{)}}{insertion}{ega25_25_182-v0010}) we therefore see that the natural right actions of
$\underline{O}(n)_S$ on the first member of (**), and the projection (**), make
$\underline{Isom}(\underline{E}_0,\underline{E})$ a \uline{principal homogeneous bundle over} $\underline{Q}(\underline{E})^*$ , \uline{with group} $\underline{O}(n)_{\underline{Q}(\underline{E})^*}$ ;
more particularly, \var{\deleted{\var{\xout{$Gl(\underline{E}_{\hand{\unsure{o|0}}})=Gl$}}{}{typewriter strike-out}{ega25_25_182-v0013}}}{}{deletion}{ega25_25_182-v0012} $Gl(n)$ is a principal homogeneous
bundle over $\underline{Q}(n)^* = \underline{Q}(\underline{Z}^n)^*$, with structure group $\underline{O}(n)_{\underline{Q}(n)^*}$ . It
follows \var{\deleted{\var{\xout{that}}{}{typewriter strike-out}{ega25_25_182-v0015}}}{}{deletion}{ega25_25_182-v0014} first that (with the notation of the proof of 16.9)
P is in fact a principal homogeneous bundle under the group $\underline{O}(n)_S$ , associated
\var{\xout{canonically}}{}{typewriter strike-out}{ega25_25_182-v0016} moreover canonically with the form $Q$ (functorially
for isomorphisms of forms, and compatibly with base chan\_%

\mspage{\hand{182}}
\tnote{25-182-1}{Accepted retains the ç/q uncertainty in reçue; the lowered doubtful comma remains adjacent to its translated word. The two parts of we retain the n/j and o/e doubts in nous. In the canceled clause, ent translates the doubtful French ending us in congrus; the doubtful le/la both translate as the. The l/i-to-l retouch in les has no English effect on the; the old \unsure{l|i} is retained here. The literal underscore at the final word break is retained in the English split chan-/ges.}
